\par\bigskip
\begin{center}
\rule{.88\textwidth}{.3pt}\par
Göttingische gelehrte Anzeigen.\quad 24. Februar 1823.\par
\rule{.88\textwidth}{.3pt}
\end{center}
\pdfbookmark[1]{24. Februar 1823}{anzeige-1823-02-24}

Eine am 2. Februar der Königlichen Societät von Herrn Hofrat \textsc{Gauss} überreichte Vorlesung mit dem Titel
\begin{center}
\textit{Theoria combinationis observationum erroribus minimis obnoxiae, pars posterior,}
\end{center}
\noindent steht in unmittelbarem Zusammenhang mit einer früheren, von der in diesen Blättern [26. Februar 1821] eine Anzeige gegeben wurde. Wir erinnern daran nur kurz: Ihr Zweck war, die sogenannte Methode der kleinsten Quadrate auf neue Weise zu begründen. Auf dieser Grundlage erscheint die Methode als die zweckmäßigste Art, Beobachtungen zu kombinieren: nicht näherungsweise, sondern mit mathematischer Schärfe; nicht nur beschränkt auf den Fall einer sehr großen Anzahl von Beobachtungen und nicht abhängig von einem hypothetischen Gesetz für die Wahrscheinlichkeit der Beobachtungsfehler, sondern in vollkommener Allgemeinheit. Der gegenwärtige zweite Teil der Untersuchung führt diese Lehre nun in einer Reihe eng damit verbundener Lehrsätze und Aufgaben weiter aus. Es entspräche der Einrich\SourceLeaf{106}%
tung dieser Blätter nicht, diesen Untersuchungen hier Schritt für Schritt zu folgen; auch ist das unnötig, da die Abhandlung selbst bereits im Druck ist. Wir begnügen uns daher damit, die Gegenstände einiger Untersuchungen anzuführen, die sich leichter für sich herausheben lassen.

Die Werte der unbekannten Größen, die der Methode der kleinsten Quadrate entsprechen und die man die \textit{sichersten Werte} nennen kann, werden durch eine bestimmte Elimination gefunden; die diesen Bestimmungen beizulegenden Gewichte durch eine unbestimmte Elimination, wie es bereits aus der \textit{Theoria motus Corporum Coelestium} bekannt ist. Hier wird auf neue Weise \textit{a priori} bewiesen, dass diese Elimination unter den gegebenen Voraussetzungen stets möglich ist. Zugleich wird eine bemerkenswerte Symmetrie unter den Koeffizienten nachgewiesen, die bei der unbestimmten Elimination hervorgehen.

So leicht und übersichtlich diese Eliminationsgeschäfte im Allgemeinen sind, so lässt sich doch nicht leugnen, dass ihre tatsächliche numerische Ausführung bei einer beträchtlichen Anzahl unbekannter Größen beschwerlich wird. Was die \textit{bestimmte} Elimination betrifft, die zur Ermittlung der sichersten Werte der unbekannten Größen ausreicht, so hat der Verfasser ein Verfahren, durch das sich die wirkliche Rechnung so weit verkürzen lässt, wie es die Natur der Sache zulässt, bereits in der \textit{Theoria Motus Corporum Coelestium} angedeutet und in einer Abhandlung im ersten Band der \textit{Commentt. Rec. Soc. R. Gott.} mit dem Titel \textit{Disquisitio de elementis ellipticis Palladis} ausführlich entwickelt. Dieses Verfahren hat zugleich den Vorteil, dass sich das Gewicht der Bestimmung der einen unbekannten Größe, die man im Rechengang zuletzt behandelt hat, von selbst ergibt. Da die Reihenfolge der unbekannten Größen völlig willkürlich ist und man also jede beliebige als letzte behandeln kann, reicht dieses Verfahren in allen Fällen aus, in denen nur das Gewicht einer einzigen unbekannten Größe verlangt wird; die beschwerliche unbestimmte Elimination wird dann vermieden.

Die inzwischen unter rechnenden Astronomen so allgemein gewordene Gewohnheit, die Methode der kleinsten Quadrate auf schwierige astronomische Rechnungen anzuwenden, etwa auf die vollständige Bestimmung von Kometenbahnen, bei denen die Zahl der unbekannten Größen bis auf sechs steigt, hat jedoch den Bedarf deutlich gemacht, die Gewichte der sichersten Werte \textit{aller} unbekannten Größen auf \SourceLeaf{107}%
eine bequemere Art als durch die unbestimmte Elimination zu finden.
Da die Bemühungen einiger Geometer\footnote{Zum Beispiel die von \textsc{Plana}. Siehe \textit{Zeitschrift für Astronomie und verwandte Wissenschaften}, Bd. 6, S. 258.} keinen Erfolg hatten, half man sich nur dadurch, dass man den oben erwähnten Algorithmus so oft ausführte, wie es unbekannte Größen gab, und dabei jedes Mal deren Reihenfolge änderte, sodass jede einmal den letzten Platz erhielt. Es scheint uns jedoch, dass dieses kunstlose Verfahren im Vergleich zur unbestimmten Elimination hinsichtlich der Kürze der Rechnung nichts gewinnt. Der Verfasser hat diesen wichtigen Gegenstand daher eigens untersucht und einen neuen Algorithmus zur Bestimmung der Gewichte der Werte \textit{aller} unbekannten Größen mitgeteilt, der die ganze Geschmeidigkeit und Kürze zu besitzen scheint, die der Gegenstand seiner Natur nach zulässt.

Der sicherste Wert einer Größe, die eine gegebene Funktion der unbekannten Größen der Aufgabe ist, wird gefunden, indem man für diese die durch die Methode der kleinsten Quadrate erhaltenen sichersten Werte einsetzt. Eine bisher nicht behandelte Aufgabe ist jedoch, wie das Gewicht zu bestimmen sei, das dieser Bestimmung beizulegen ist. Die hier gegebene Lösung verdient umso mehr die Aufmerksamkeit der rechnenden Astronomen, als sich zeigt, dass mehrere von ihnen dabei früher nicht richtig vorgegangen sind.

Bekanntlich bildet die Summe der Quadrate der Unterschiede zwischen den unmittelbar beobachteten Größen und denjenigen Werten, die sich für sie als Funktionen der unbekannten Größen durch Einsetzen der sichersten Werte der Letzteren ergeben, ein absolutes Minimum. (Sind die Beobachtungen unterschiedlich zuverlässig, müssen die Quadrate vor dem Addieren zunächst mit den jeweiligen Gewichten multipliziert werden.) Sobald man daher einer unbekannten Größe einen Wert zuweist, der von ihrem sichersten Wert abweicht, fällt ein ähnliches Aggregat stets größer aus als das erwähnte Minimum, gleichgültig, wie man die übrigen unbekannten Größen bestimmt. Die übrigen unbekannten Größen lassen sich jedoch nur auf eine einzige Weise so bestimmen, dass die Vergrößerung des Aggregats möglichst klein ist oder das Aggregat selbst zu einem relativen Minimum wird. Diese vom Ver\SourceLeaf{108}%
fasser hier durchgeführte Untersuchung führt zu einigen interessanten Wahrheiten, die die ganze Lehre in noch vielseitigeres Licht rücken.
Es kommt zuweilen vor, dass man erst dann, wenn man eine umfangreiche Rechnung über eine Reihe von Beobachtungen vollständig durchgeführt hat, von einer neuen Beobachtung erfährt, die man gern noch einbezogen hätte. In vielen Fällen kann es wünschenswert sein, deshalb nicht die gesamte Eliminationsarbeit von vorn beginnen zu müssen, sondern die Änderung bestimmen zu können, die durch die Hinzunahme der neuen Beobachtung bei den sichersten Werten und ihren Gewichten entsteht. Der Verfasser hat daher diese Aufgabe eigens behandelt, ebenso wie die verwandte Aufgabe, bei der man sich nachträglich veranlasst sieht, einer bereits verwendeten Beobachtung ein anderes als das ihr zugewiesene Gewicht zu geben und die Änderungen der Endergebnisse zu erhalten, ohne die Rechnung von vorn zu wiederholen.

In einer eigenen Abhandlung in der \textit{Zeitschrift für Astronomie und verwandte Wissenschaften} [März und April 1816] hatte der Verfasser gezeigt, wie sich der \textit{wahrscheinliche} Fehler einer Beobachtungsgattung – bislang das gebräuchliche Maß ihrer Unsicherheit – aus einer hinreichenden Anzahl wirklicher Beobachtungsfehler näherungsweise finden lässt. Dieses Verfahren und überhaupt der Gebrauch des wahrscheinlichen Fehlers hängen jedoch von der hypothetischen Form der Wahrscheinlichkeit der einzelnen Fehler ab, und mussten davon abhängen. Im ersten Teil der gegenwärtigen Abhandlung wird zwar gezeigt, wie sich aus denselben Daten näherungsweise der mittlere Fehler der Beobachtungen – als zweckmäßiges Maß ihrer Ungenauigkeit – bestimmen lässt. Dennoch bleibt die Schwierigkeit, dass man nur selten oder fast nie mit aller Genauigkeit die wahre Größe einer Anzahl wirklicher Beobachtungsfehler kennen kann. In der Praxis hat man daher stets die Unterschiede zwischen den Ergebnissen der Beobachtungen und den Resultaten der Rechnung zugrunde gelegt, die sich aus den durch die Methode der kleinsten Quadrate ermittelten sichersten Werten der unbekannten Größen ergeben, von denen die Beobachtungen abhängen. Da man aber nicht berechtigt ist, die sichersten Werte für die wahren Werte selbst zu halten, erkennt man leicht, dass dieses Verfahren den wahrscheinlichen wie den mittleren Fehler stets \textit{zu klein} ergeben muss und den Beobachtungen sowie den daraus gezo\SourceLeaf{109}%
genen Ergebnissen daher eine größere Genauigkeit zuschreibt, als sie tatsächlich besitzen. Freilich fällt diese Ungenauigkeit wenig ins Gewicht, wenn die Zahl der Beobachtungen die Zahl der unbekannten Größen um ein Vielfaches übersteigt. Doch verlangt einerseits die Würde der Wissenschaft, vollständig und genau zu überblicken, wie viel man auf diese Weise fehlzugehen riskiert; andererseits sind nach diesem fehlerhaften Verfahren tatsächlich schon oft Rechenergebnisse in wichtigen Fällen aufgestellt worden, in denen jene Voraussetzung nicht erfüllt war. Der Verfasser hat diesen Gegenstand daher eigens untersucht und ein höchst bemerkenswertes, sehr einfaches Ergebnis erhalten. Um den nach dem genannten fehlerhaften Verfahren gefundenen mittleren Fehler in den richtigen umzuwandeln, muss man ihn lediglich mit
\[
\surd\frac{\pi-\rho}{\pi}
\]
multiplizieren, wobei \(\pi\) die Anzahl der Beobachtungen und \(\rho\) die Anzahl der unbekannten Größen bezeichnet.

Die letzte Untersuchung betrifft noch die Bestimmung des Genauigkeitsgrades, der dieser Bestimmung des mittleren Fehlers selbst beizulegen ist; ihre Ergebnisse muss man jedoch in der Abhandlung selbst nachlesen.
\par\bigskip
\begin{center}
\rule{.88\textwidth}{.3pt}\par
\end{center}
